% !TeX root = main.tex
\chapter{Related Work}
\label{chap:related_work}

This chapter reviews \emph{Homogeneous Synchronous FL} (Section~\ref{sec:related_sync_fl}), \emph{Submodel FL} (Section~\ref{sec:related_model_heterogeneous_fl}), \emph{Asynchronous FL} (Section~\ref{sec:related_async_fl}), and \emph{Asynchronous Submodel FL} (Section~\ref{sec:async_submodel_related}). Appendix~\ref{chap:notation} summarises the notation used, and Appendix~\ref{appendix:algorithms} gives unified pseudocode for each implemented FL algorithm.

Other previous work includes addressing data heterogeneity through \textbf{optimisation corrections} \parencite{li2020fedprox,karimireddy2020scaffold,wang2020fednova}, and addressing systems heterogeneity through \textbf{client selection and pacing} \parencite{nishio2018fedcs,lai2021oort,chai2020tifl}; these are complementary to the FL methods evaluated here.

% All methods studied here retain the federated objective defined in Section~\ref{sec:problem_setting}. The scientific question is therefore not what objective the methods claim to solve, but how they change the \emph{systems path} by which information from heterogeneous clients reaches the server model.

\section{Homogeneous Synchronous FL}
\label{sec:related_sync_fl}

The reference point for the two following method families is standard synchronous FL. In each round, the server broadcasts \(w^{(t)}\), each selected client performs local optimisation, and the server aggregates the returned models by an example-weighted average:
\[
    \tilde{w}_k^{(t+1)} \leftarrow \operatorname{ClientUpdate}(w^{(t)}, D_k),
    \qquad
    w^{(t+1)}
    =
    \frac{\sum_{k\in\mathcal{C}_t} n_k \tilde{w}_k^{(t+1)}}
         {\sum_{h\in\mathcal{C}_t} n_h}.
\]

\subsection{\texttt{FedAvg}}
\label{sec:fedavg_related}

\texttt{FedAvg} is the synchronous FL baseline \parencite{mcmahan2017fedavg} and follows the update rule above. In each round, it trains the full model on each selected client and waits for all selected clients' responses before aggregation. In chapter~\ref{chap:evaluation}, the comparison of \texttt{FedAvg} with other methods not only demonstrates the negative effects of heterogeneity, but also shows the effects of submodel FL and asynchronous FL in addressing these issues.

\section{Submodel FL}
\label{sec:related_model_heterogeneous_fl}

\textbf{Submodel FL}, also known as \textbf{partial-training-based model-heterogeneous FL}, aims to address compute heterogeneity. It relaxes the constraint that same-size models must be trained on all clients and allows weaker clients to train smaller subnetworks. The methods studied here share a common partial-training structure: 
the server uses an \textbf{extraction policy} \( \Psi_{\mathcal{M}} \) to choose a parameter subset \( I_k^{(t)} \), the client trains the model initialised with weights \( w_{I_k^{(t)}}^{(t)} \), and the server aggregates each parameter only over clients that actually trained it.

\[
    I_k^{(t)} = \Psi_{\mathcal{M}}\!\left(w^{(t)}, r_k, t\right),
    \qquad
    w_{k,0}^{(t)} = w^{(t)}_{I_k^{(t)}}.
\]
Here, \(r_k\) is the fixed model rate assigned to client \(k\). The following equation gives the per-parameter aggregation, where \(\mathcal{U}_j^{(t)}\) is the set of clients that trained parameter \(j\) in round \(t\), and \([w]_j\) denotes the \(j\)th parameter of model \(w\):
\[
    \mathcal{U}_j^{(t)} = \{k\in\mathcal{C}_t: j\in I_k^{(t)}\},
    \qquad
    [w^{(t+1)}]_j
    =
    \frac{\sum_{k\in\mathcal{U}_j^{(t)}} n_k [\tilde{w}_k^{(t+1)}]_j}
         {\sum_{h\in\mathcal{U}_j^{(t)}} n_h},
    \quad
    \mathcal{U}_j^{(t)}\neq\emptyset.
\]

Figure~\ref{fig:model_heterogeneity_method_overview} shows the distinction between the three submodel FL methods. Across synchronous rounds, \texttt{HeteroFL} \textbf{reuses fixed nested prefixes}, \texttt{FedRolex} \textbf{rotates the selected window}, and \texttt{FIARSE} applies an \textbf{importance-aware sparse mask}.

\begin{figure}[H]
\centering
\includegraphics[width=\textwidth]{figures/model_heterogeneity_method_overview.pdf}
\caption{\textbf{Method-overview for the three submodel FL baselines.}}
\label{fig:model_heterogeneity_method_overview}
\end{figure}

\subsection{\texttt{HeteroFL}}
\label{sec:heterofl_related}

\texttt{HeteroFL} instantiates the extractor as \(\Psi_{\mathrm{nested}}\), a \textbf{static nested} prefix selector \parencite{diao2021heterofl}. For layer \(\ell\) with \(d_\ell\) structured units, the layer-wise index set is
\[
    I_{k,\ell}^{(t)}
    =
    \{1,\dots,\lceil r_k d_\ell\rceil\},
    \qquad
    I_k^{(t)}=\Psi_{\mathrm{nested}}\!\left(w^{(t)}, r_k\right).
\]
The resulting submodels are nested, so every lower-rate model is contained in every higher-rate model. Under per-parameter aggregation, prefix parameters are averaged over more clients, while outer parameters are updated only by clients with a larger model rate \(r_k\). This ties each client to a fixed parameter region which stabilises training but underexposes some parameters to weak clients. 

\subsection{\texttt{FedRolex}}
\label{sec:fedrolex_related}

\texttt{FedRolex} instantiates the extractor as \(\Psi_{\mathrm{roll}}\), a \textbf{rolling-window} selector \parencite{alam2022fedrolex}. Rather than selecting the same nested prefix, it shifts the window over the parameter space:
\[
    I_k^{(t)}=\Psi_{\mathrm{roll}}\!\left(w^{(t)}, r_k, \rho_t\right),
\]
where the offset \(\rho_t\) advances with the round index. This gives lower-rate clients opportunities to train different parameter regions over time. However, per-client submodel training can be less stable.

% \begin{definition}
% \begin{lemma}[FedRolex coverage lemma {\normalfont\parencite{alam2022fedrolex}}]
% Given \(I\) candidate indices and one uniformly chosen index per round, the expected number of rounds required to choose all indices at least once is
% \[
%     I \left(\frac{1}{I} + \frac{1}{I-1} + \cdots + \frac{1}{1}\right),
% \]
% which is asymptotically \(I \log I\) as \(I \to \infty\).
% \end{lemma}
% \end{definition}

\subsection{\texttt{FIARSE}}
\label{sec:fiarse_related}

\texttt{FIARSE} instantiates the extractor as \(\Psi_{\mathrm{importance}}\), an \textbf{importance-aware sparse mask} \parencite{wu2024fiarse}. It is motivated by previous work on sparsity and pruning \parencite{hoefler2021sparsity}. The importance score \(s_j^{(t)}\) for parameter \(j\) in round \(t\) is just its magnitude
\[
    s_j^{(t)}=|[w^{(t)}]_j|,
    \qquad
    I_k^{(t)}
    =
    \Psi_{\mathrm{importance}}\!\left(w^{(t)}, r_k\right)
    =
    \{j:s_j^{(t)}\geq\theta_k^{(t)}\},
\]
\(\theta_k^{(t)}\) is the threshold that keeps the top \(r_k\) fraction of parameters. Therefore, each submodel is expected to be a high-quality sparse approximation of the global model, hence improving submodel quality. One caveat is that \cite{wu2024fiarse}'s efficiency argument relies on sparse execution, while the evaluation here runs the masks through dense PyTorch kernels on Raspberry Pi CPUs, separating submodel-quality gains from sparse-kernel runtime benefits.


\section{Asynchronous FL}
\label{sec:related_async_fl}

Asynchronous FL aims to address straggler-induced heterogeneity. It removes the restriction in synchronous FL that waits for the entire client set to finish, which can be costly when device capabilities or network conditions differ across clients. Using the staleness notation from Section~\ref{sec:background_straggler_heterogeneity}, let server step \(s\) index server updates. If update \(i\) was trained from server version \(v_i\), then its arrival staleness is
\[
    \tau_i = s - v_i.
\]

\begin{figure}[H]
\centering
\begin{adjustbox}{max width=\textwidth}
\begin{tikzpicture}[
  >=Latex,
  font=\scriptsize,
  outer/.style={draw=black!18, fill=gray!3, rounded corners=5pt, minimum width=15.1cm, minimum height=6.5cm},
  panel/.style={draw=black!18, fill=white, rounded corners=4pt, minimum width=4.15cm, minimum height=2.05cm},
  clientfast/.style={draw=darkblue!70!black, fill=darkblue!7, rounded corners=3pt, align=center, minimum width=1.35cm, minimum height=0.55cm},
  clientslow/.style={draw=paletteOrange!85!black, fill=paletteOrange!12, rounded corners=3pt, align=center, minimum width=1.35cm, minimum height=0.55cm},
  updatefast/.style={draw=darkblue!70!black, fill=darkblue!8, circle, minimum size=6.7mm, inner sep=0pt, font=\scriptsize\bfseries},
  updateslow/.style={draw=paletteOrange!85!black, fill=paletteOrange!13, circle, minimum size=6.7mm, inner sep=0pt, font=\scriptsize\bfseries},
  box/.style={draw=black!28, fill=white, rounded corners=3pt, align=center, inner sep=4pt},
  arrow/.style={->, line width=0.65pt, draw=black!58},
  mutedarrow/.style={->, line width=0.55pt, draw=black!38},
  methodarrow/.style={->, line width=0.8pt},
  note/.style={align=center, text=black!62},
  barlabel/.style={font=\tiny, text=black!58, align=center}
]
\node[outer] at (0,0) {};
\node[font=\bfseries\small, text=black!68] at (0,2.67) {Asynchronous execution choices};

\node[box, text width=1.25cm] (server) at (-6.65,1.25) {server\\version \(v\)};
\node[clientfast] (fast) at (-4.45,1.70) {fast\\clients};
\node[clientslow] (slow) at (-4.45,0.78) {slow\\clients};
\draw[mutedarrow] (server.east) -- (fast.west);
\draw[mutedarrow] (server.east) -- (slow.west);
\node[box, minimum width=1.75cm, minimum height=1.05cm] (buffer) at (0.85,1.22) {};
\node[updatefast] (uone) at (-2.70,1.82) {\(u_1\)};
\node[updatefast] (utwo) at (-1.86,1.20) {\(u_2\)};
\node[updateslow] (uthree) at (-1.00,0.58) {\(u_3\)};
\node[barlabel] at (-2.70,2.3) {\(\tau=0\)};
\node[barlabel] at (-1.86,0.79) {\(\tau=1\)};
\node[barlabel, text=paletteOrange!90!black] at (-1.00,0.16) {\(\tau=3\)};
\draw[mutedarrow] (fast.east) -- (uone.west);
\draw[mutedarrow] (fast.east) -- (utwo.west);
\draw[mutedarrow] (slow.east) -- (uthree.west);

\node[font=\scriptsize\bfseries, text=black!68] at (0.8,2) {buffer fills at \(K=3\)};
\draw[draw=darkblue!70!black, fill=darkblue!8, line width=0.55pt] (0.21,0.9) rectangle (0.55,1.5);
\draw[draw=darkblue!70!black, fill=darkblue!8, line width=0.55pt] (0.67,0.9) rectangle (1.01,1.5);
\draw[draw=paletteOrange!85!black, fill=paletteOrange!13, line width=0.55pt] (1.13,0.9) rectangle (1.47,1.5);
\node[barlabel, text=darkblue!80!black] at (0.85,0.50) {no buffer for \texttt{FedAsync}};
\draw[arrow] (uone.east) -- (buffer.west);
\draw[arrow] (utwo.east) -- (buffer.west);
\draw[arrow] (uthree.east) -- (buffer.west);

\node[box, text width=1.55cm, minimum height=0.78cm] (aggregate) at (3.65,1.22) {weighted\\aggregate};
\node[box, text width=1.35cm, minimum height=0.78cm] (nextmodel) at (6.35,1.22) {next model\\\(w^{[s+1]}\)};
\draw[arrow] (buffer.east) -- (aggregate.west);
\draw[arrow] (aggregate.east) -- (nextmodel.west);

\draw[draw=black!18, line width=0.45pt] (-6.35,-0.1) -- (6.35,-0.1);

\node[panel] (fedasyncpanel) at (-4.85,-1.95) {};
\node[panel] (fedbuffpanel) at (0,-1.95) {};
\node[panel] (fedstalepanel) at (4.85,-1.95) {};

\node[font=\bfseries\scriptsize, text=darkblue] at (-4.85,-0.6) {\texttt{FedAsync}: immediate update};
\node[font=\scriptsize, text=black!62] at (-4.85,-1.2) {\(\lambda_i=\lambda_0\varphi(\tau_i)\)};
\node[updatefast, minimum size=6.2mm] (asyncupdate) at (-6.50,-2.15) {\(u_i\)};
\node[box, text width=0.85cm, minimum height=0.55cm] (asyncmix) at (-5.1,-2.15) {update\\now};
\node[box, text width=0.85cm, minimum height=0.55cm] (asyncnext) at (-3.45,-2.15) {next\\model};
\draw[arrow] (asyncupdate.east) -- (asyncmix.west);
\draw[arrow] (asyncmix.east) -- (asyncnext.west);
\node[barlabel] at (-6.50,-2.63) {\(\tau_i\)};
\node[font=\tiny\bfseries, text=black!55, align=center] at (-5.0,-2.82) {no buffer};

\node[font=\bfseries\scriptsize, text=fedbuffteal] at (0,-0.6) {\texttt{FedBuff}: stale decay};
\node[font=\scriptsize, text=black!62] at (0,-1.2) {\(q_i=(1+\tau_i)^{-1/2}\)};
\draw[->, draw=black!35, line width=0.45pt] (-1.65,-2.67) -- (1.65,-2.67);
\draw[->, draw=black!35, line width=0.45pt] (-1.45,-2.90) -- (-1.45,-1.62);
\draw[fill=darkblue!18, draw=darkblue!75!black] (-1.20,-2.67) rectangle (-0.78,-1.84);
\draw[fill=darkblue!12, draw=darkblue!75!black] (-0.35,-2.67) rectangle (0.07,-2.08);
\draw[fill=paletteOrange!12, draw=paletteOrange!90!black] (0.50,-2.67) rectangle (0.92,-2.39);
\node[barlabel] at (-0.99,-2.84) {\(\tau=0\)};
\node[barlabel] at (-0.14,-2.84) {\(\tau=1\)};
\node[barlabel, text=paletteOrange!90!black] at (0.71,-2.84) {\(\tau=3\)};
\node[barlabel, rotate=90] at (-1.72,-2.20) {weight \(\alpha\)};
\node[font=\tiny\bfseries, text=paletteOrange!90!black, align=center] at (1.28,-2.04) {stale\\low};

\node[font=\bfseries\scriptsize, text=fedstalebrown] at (4.85,-0.6) {\texttt{FedStaleWeight}: stale weights};
\node[font=\scriptsize, text=black!62] at (4.85,-1.2) {\(q_i=\tau_iK+1\)};
\draw[->, draw=black!35, line width=0.45pt] (3.20,-2.67) -- (6.50,-2.67);
\draw[->, draw=black!35, line width=0.45pt] (3.40,-2.90) -- (3.40,-1.62);
\draw[fill=darkblue!10, draw=darkblue!75!black] (3.65,-2.67) rectangle (4.07,-2.46);
\draw[fill=darkblue!12, draw=darkblue!75!black] (4.50,-2.67) rectangle (4.92,-2.20);
\draw[fill=paletteOrange!18, draw=paletteOrange!90!black] (5.35,-2.67) rectangle (5.77,-1.84);
\node[barlabel] at (3.86,-2.84) {\(\tau=0\)};
\node[barlabel] at (4.71,-2.84) {\(\tau=1\)};
\node[barlabel, text=paletteOrange!90!black] at (5.56,-2.84) {\(\tau=3\)};
\node[barlabel, rotate=90] at (3.13,-2.20) {weight \(\alpha\)};
\node[font=\tiny\bfseries, text=paletteOrange!90!black, align=center] at (6.18,-2.04) {stale\\high};
\end{tikzpicture}
\end{adjustbox}
\caption{\textbf{Method-overview for the three asynchronous FL baselines.}}
\label{fig:async_method_overview}
\end{figure}

Figure~\ref{fig:async_method_overview} shows the distinction between the three asynchronous FL methods. \texttt{FedAsync} applies each reply immediately to update the global model, whereas \texttt{FedBuff} and \texttt{FedStaleWeight} first collect replies in a buffer of size \(K\), then aggregate the replies with stale-decay and stale-weight normalisation respectively.



\subsection{\texttt{FedAsync}}
\label{sec:fedasync_related}

\texttt{FedAsync} is the fully asynchronous baseline: every received client model immediately updates the global model, with no aggregation buffer \parencite{xie2019asynchronous}. Each arrival applies a weighted update between the current model \(w^{[s]}\) and the returned model \(\tilde{w}_i\),
\[
    w^{[s+1]}
    =
    (1-\lambda_i) w^{[s]} + \lambda_i \tilde{w}_i,
    \qquad
    \lambda_i = \lambda_0 \varphi(\tau_i).
\]
Here \(\lambda_i\) is the immediate-arrival update coefficient, \(\lambda_0\) is the base rate, and \(\varphi(\tau_i)\) downweights stale replies. \cite{xie2019asynchronous} considers constant weighting \(\varphi(\tau)=1\), polynomial decay \(\varphi(\tau)=(1 + \tau)^{-a}\), and hinge decay, which keeps full weight up to a delay threshold and then decays.

\subsection{\texttt{FedBuff}}
\label{sec:fedbuff_related}

\texttt{FedBuff} sits between synchronous \texttt{FedAvg} and one-update-at-a-time \texttt{FedAsync} \parencite{nguyen2022fedbuff}. The server keeps several client jobs in flight. As clients finish, their updates are appended to a buffer; when \(K\) replies have accumulated, the server normalises their weights, applies one aggregate update, clears the buffer, and dispatches more work. This enables non-blocking execution while avoiding a separate global update for every single reply arrival.

For a filled buffer \(\mathcal{B}_s\), the buffered update is
\[
    \alpha_i = \frac{q_i}{\sum_{u\in\mathcal{B}_s} q_u},
    \qquad
    w^{[s+1]}
    =
    w^{[s]} - \eta_s
    \sum_{i\in\mathcal{B}_s}\alpha_i\Delta_i .
\]
Here \(\Delta_i\) is the model delta from reply \(i\). \(q_i\) is the staleness score assigned to reply \(i\), and \(\alpha_i\) is the normalised buffer weight. \texttt{FedBuff} assigns each reply \(i\) the \textbf{staleness score}:
\[
    q_i = (1+\tau_i)^{-1/2}.
\]
Older replies therefore receive smaller buffer weights, but still contribute as part of the aggregate. The buffer size \(K\) controls how many updates are combined per server step, while client concurrency controls how many clients may train at once. This reduces straggler blocking, supports secure aggregation more naturally than \texttt{FedAsync}, and makes \(K\) the evaluation variable for speed, update diversity, and slow-device representation.

% \begin{definition}
% \begin{theorem}[FedBuff ergodic convergence {\normalfont\parencite{nguyen2022fedbuff}}]
% Assume unbiased stochastic gradients, bounded local and global variance, bounded gradients, \(L\)-smooth local objectives, and bounded staleness. Let
% \[
%     A(Q) := \sum_{q=0}^{Q-1} \eta_\ell^{(q)},
%     \qquad
%     B(Q) := \sum_{q=0}^{Q-1} \left(\eta_\ell^{(q)}\right)^2.
% \]
% If the local and server learning rates satisfy \(\eta_s \eta_\ell^{(q)} Q \leq 1/L\) for all local steps \(q\), then the buffered asynchronous iterates satisfy the ergodic bound
% \[
% \begin{aligned}
%     \frac{1}{T} \sum_{t=0}^{T-1} \left\lVert \nabla f(w^t) \right\rVert^2
%     &\leq \frac{2 \left(f(w^0) - f^* \right)}{\eta_s A(Q) T}
%     + \frac{L}{2}\frac{\eta_s B(Q)}{A(Q)} \sigma_\ell^2 \\
%     &\quad + 3L^2 Q B(Q)\left(\eta_s^2 \tau_{\max,K}^2 + 1\right)
%     \left(\sigma_\ell^2 + \sigma_g^2 + G\right).
% \end{aligned}
% \]
% \end{theorem}
% \end{definition}

% The original analysis gives convergence under standard smooth non-convex assumptions and bounded staleness; here, the key point is that \texttt{FedBuff} exposes buffer size and staleness as optimisation variables rather than implementation details.

\subsection{\texttt{FedStaleWeight}}
\label{sec:fedstaleweight_related}

\texttt{FedStaleWeight} keeps the same buffered execution model as \texttt{FedBuff}, but changes the staleness score and hence the buffer weights \parencite{ma2024fedstaleweight}. Fast clients can dominate asynchronous training because they complete more trips; this creates a problem when data are non-IID, especially when device speed and data distribution are correlated. Therefore, \texttt{FedStaleWeight} upweights stale replies to counteract this bias. The \textbf{staleness score} for reply \(i\) is:
\[
    q_i = \tau_iK + 1.
\]
Staler replies therefore receive larger weights after normalisation. This is useful when staleness tracks under-represented client groups, but hurts performance when clients' data, and hence their update gradients, are IID. The bias correction effect is evaluated in Section~\ref{sec:eval_device_correlated_skew}.

% \begin{definition}
% \begin{theorem}[FedStaleWeight ergodic convergence {\normalfont\parencite{ma2024fedstaleweight}}]
% Under the same smoothness, bounded-variance, bounded-gradient, and bounded-staleness assumptions used for buffered asynchronous FL, let
% \[
%     A(Q) := \sum_{q=0}^{Q-1} \eta_\ell^{(q)},
%     \qquad
%     B(Q) := \sum_{q=0}^{Q-1} \left(\eta_\ell^{(q)}\right)^2,
% \]
% and let \(U(\tau_{\max,b})\) denote the maximum squared deviation between equal buffer weighting and the normalised stale-weighted coefficients. If \(\eta_s \eta_\ell^{(q)} Q \leq 1/L\) for all local steps \(q\), then the FedStaleWeight iterates satisfy
% \[
% \begin{aligned}
%     \frac{1}{T}\sum_{t=0}^{T}\left\lVert \nabla f(w^t) \right\rVert^2
%     &\leq \frac{2\left(f(w^0)-f(w^*)\right)}{\eta_s A(Q) T} \\
%     &\quad + 12 \eta_s^2 \tau_{\max,b}^2 B(Q) Q L^2
%       \left(\frac{b\tau_{\max,b}+1}{\tau_{\max,b}+1}\right)^2
%       \left(\sigma_\ell^2 + \sigma_g^2 + G\right) \\
%     &\quad + 12 B(Q) Q L^2 \left(\sigma_\ell^2 + \sigma_g^2 + G\right)
%       + 4b^2 U(\tau_{\max,b})G.
% \end{aligned}
% \]
% \end{theorem}
% \end{definition}

% The original analysis places this stale-weighted update rule in the same smooth non-convex framework as buffered asynchronous FL, with an additional term capturing the deviation induced by reweighting. This separates the systems question of slow-client representation from the optimisation question of whether stale-aware rebalancing remains analysable.



\section{Asynchronous Submodel FL}
\label{sec:async_submodel_related}


\paragraph{Gap between method families.}
Submodel FL and asynchronous FL address compute and straggler-induced heterogeneity respectively. Submodel FL allows reduced-size models for weaker clients, while asynchronous FL relaxes waiting requirements. In a realistic edge FL setting these heterogeneities can coincide: a weak client may train a reduced submodel and still return a stale update. \cite{zhang2023timelyfl} proposes \texttt{TimelyFL}, using client workload coordination to make asynchronous participation timely. The gap left for this dissertation is the complementary server-side question: once a buffer contains stale partial submodel updates, how should the server aggregate parameters with unequal and dynamic coverage?

\paragraph{Why coverage becomes dynamic.}
Parameter coverage is already central to submodel FL. In synchronous \texttt{HeteroFL}-style aggregation, a parameter is updated only by clients whose assigned submodel contains it, so coverage is tied to the selected client set and model-rate assignment \parencite{diao2021heterofl}. \texttt{FedRolex} changes the selected window over rounds to reduce persistent concentration in one parameter region \parencite{alam2022fedrolex}. In asynchronous submodel FL, whichever clients finish first determine both the normalised buffer weights and the submodel regions present in that server step. For a fixed asynchronous server step \(s\), with accepted buffer \(\mathcal{B}_s\), the coverage mass for parameter \(j\) is
\[
    c_j^{[s]}
    =
    \sum_{i\in\mathcal{B}_s:\, j\in I_i}\alpha_i,
    \qquad
    \alpha_i
    =
    \frac{q_i}{\sum_{u\in\mathcal{B}_s}q_u}.
\]
Here, \(\alpha_i\) is the normalised buffer weight and \(q_i\) is the staleness score for update \(i\). This coverage mass \(c_j^{[s]}\) is dynamic, and can differ across parameters within the same server step.

Figure~\ref{fig:async_submodel_coverage_mismatch} illustrates the resulting aggregation problem: low-coverage parameter blocks contribute a smaller component to the buffered update even when their local delta is useful. This can misalign the coverage-scaled step with the full-coverage step, which can hurt convergence and final performance.

\begin{figure}[H]
\centering
\begin{adjustbox}{max width=\textwidth}
\begin{tikzpicture}[
  >=Latex,
  font=\small,
  panel/.style={draw=black!18, fill=gray!2, rounded corners=4pt, minimum width=7.60cm, minimum height=6.35cm},
  title/.style={font=\small\bfseries, text=black!72},
  smalllabel/.style={font=\footnotesize, align=center, text=black!67},
  rowlabel/.style={font=\footnotesize, align=center, text=black!65},
  blocklabel/.style={font=\footnotesize\bfseries, align=center},
  vec/.style={-{Latex[length=2.8mm]}, line width=0.95pt},
  resultant/.style={-{Latex[length=3.1mm]}, line width=1.15pt}
]
\node[panel] at (-4.08,0) {};
\node[panel] at (4.08,0) {};

\node[title] at (-4.08,2.78) {Parameter-block coverage};
\node[smalllabel] at (-4.08,2.38) {nested submodels cover prefix blocks more};

\draw[draw=black!45, line width=0.65pt] (-6.55,-1.62) rectangle (-1.75,1.62);
\draw[draw=paletteGreen!75!black, fill=paletteGreen!24, line width=0.65pt] (-6.55,-1.62) rectangle (-4.95,1.62);
\draw[draw=darkblue!75!black, fill=darkblue!14, line width=0.65pt] (-4.95,-1.62) rectangle (-3.35,1.62);
\draw[draw=paletteOrange!90!black, fill=paletteOrange!18, line width=0.65pt] (-3.35,-1.62) rectangle (-1.75,1.62);

\node[blocklabel, text=paletteGreen!55!black] at (-5.75,0.52) {prefix\\block};
\node[blocklabel, text=darkblue!75!black] at (-4.15,0.52) {middle\\block};
\node[blocklabel, text=paletteOrange!90!black] at (-2.55,0.52) {outer\\block};
\node[smalllabel, text=paletteGreen!55!black] at (-5.75,-0.78) {\(c=1.00\)};
\node[smalllabel, text=darkblue!75!black] at (-4.15,-0.78) {\(c=0.65\)};
\node[smalllabel, text=paletteOrange!90!black] at (-2.55,-0.78) {\(c=0.25\)};
\node[smalllabel] at (-4.08,-2.36) {coverage mass varies across the full model};

\node[title] at (4.08,2.78) {Buffered update components};
\node[smalllabel] at (4.08,2.38) {same local deltas, different coverage scaling};

\node[rowlabel] at (1.64,1.12) {full\\coverage\\\((1,1,1)\)};
\coordinate (eq0) at (2.70,1.16);
\coordinate (eq1) at (4.35,1.16);
\coordinate (eq2) at (5.60,1.81);
\coordinate (eq3) at (6.85,1.16);
\draw[vec, draw=paletteGreen!75!black] (eq0) -- (eq1);
\draw[vec, draw=darkblue!75!black] (eq1) -- (eq2);
\draw[vec, draw=paletteOrange!90!black] (eq2) -- (eq3);
\fill[black!68] (eq0) circle (1.7pt);
\draw[resultant, draw=black!64] (2.70,0.34) -- (6.85,0.34);
\node[smalllabel, text=black!58] at (4.78,0.03) {full-coverage step};

\node[rowlabel] at (1.64,-1.32) {partial\\coverage\\\((1,.65,.25)\)};
\coordinate (un0) at (2.70,-1.28);
\coordinate (un1) at (4.35,-1.28);
\coordinate (un2) at (5.16,-0.86);
\coordinate (un3) at (5.48,-1.02);
\draw[vec, draw=paletteGreen!75!black] (un0) -- (un1);
\draw[vec, draw=darkblue!75!black] (un1) -- (un2);
\draw[vec, draw=paletteOrange!90!black] (un2) -- (un3);
\fill[black!68] (un0) circle (1.7pt);
\draw[resultant, draw=black!64] (2.70,-1.98) -- (5.48,-1.72);
\node[smalllabel, text=black!58] at (4.09,-2.30) {coverage-scaled step};
\end{tikzpicture}
\end{adjustbox}
\caption{\textbf{Coverage mismatch in asynchronous submodel aggregation.} Left: nested submodels give different parameter blocks different coverage mass \(c_j^{[s]}\). Right: partial coverage scales update components unevenly, so the coverage-scaled step can deviate.}
\label{fig:async_submodel_coverage_mismatch}
\end{figure}

\paragraph{\texttt{FedCover} motivation.}
\label{sec:coverage_mismatch_related}
\label{sec:fedcover_related}
A direct baseline is therefore \texttt{Async-HeteroFL}, which is \texttt{FedBuff}-style buffering combined with \texttt{HeteroFL}-style nested slicing. This baseline accepts stale partial replies, but applies the buffered aggregate without correcting for parameters whose coverage mass \(c_j^{[s]}\) is small, producing the coverage-scaled step. \texttt{FedCover} addresses this server-side aggregation problem by applying a capped inverse-coverage correction to under-covered parameter blocks. Section~\ref{sec:fedcover_impl} defines the algorithm, and Section~\ref{sec:eval_async_submodel_fl} evaluates whether this correction improves the performance--runtime trade-off under various conditions.


% \section{From Methods to Evaluation Design}
% \label{sec:related_method_evaluation_bridge}

% The methods above are selected because they exercise different deployed-system failure modes, not because they exhaust the FL literature. Table~\ref{tab:method_evaluation_bridge} summarises how each family maps onto the dissertation's evaluation questions. The common requirement is that all methods must run through the same physical cluster, configuration surface, logging path, and result-collection workflow; otherwise differences in quality or runtime could be artifacts of the harness rather than consequences of the method.

% \begin{table}[!htbp]
% \centering
% \caption{Method families and the evaluation evidence they require.}
% \label{tab:method_evaluation_bridge}
% \footnotesize
% \begin{tabularx}{\textwidth}{>{\raggedright\arraybackslash}p{0.24\textwidth}>{\raggedright\arraybackslash}p{0.20\textwidth}>{\raggedright\arraybackslash}X>{\centering\arraybackslash}p{0.12\textwidth}}
% \toprule
% Method family & Algorithms & Evaluation role & Questions \\
% \midrule
% Full-model synchronous baseline & \texttt{FedAvg} & Reference point for full-model quality, synchronous runtime, and client-trip cost across both axes & C1, S1 \\
% Submodel FL methods & \texttt{HeteroFL}, \texttt{FedRolex}, \texttt{FIARSE} & Client-capacity comparison using submodel quality, train-round latency, and device/rate timing & C1--C3 \\
% Immediate asynchronous method & \texttt{FedAsync} & Straggler-induced heterogeneity comparison using time-to-target, slow-device update share, and stale-update behaviour & S1--S2 \\
% Buffered asynchronous methods & \texttt{FedBuff}, \texttt{FedStaleWeight} & Straggler-induced heterogeneity and representation comparison using staleness, buffer behaviour, and aggregate weight share & S1--S3 \\
% \bottomrule
% \end{tabularx}
% \end{table}

% Chapter~\ref{chap:implementation} describes the shared deployment and measurement controls that make these method families comparable on the same physical cluster.
